\documentclass[tikz,border=10pt]{standalone}
\usepackage{tikz}
\usepackage{amsmath}
\usepackage{amssymb}
\usepackage{ctex}
\usetikzlibrary{calc, positioning, arrows.meta, decorations.pathreplacing}

\begin{document}
\begin{tikzpicture}[>=Stealth,
  datanode/.style={
    draw=gray!60, fill=white, rounded corners=6pt,
    minimum width=1.9cm, minimum height=1.1cm,
    inner sep=5pt, align=center, font=\small
  },
  opnode/.style={
    draw=#1!70, fill=#1!12, rounded corners=6pt,
    minimum width=1.9cm, minimum height=1.1cm,
    inner sep=5pt, align=center, font=\small
  },
  lossnode/.style={
    draw=red!70, fill=red!15, rounded corners=6pt,
    minimum width=1.9cm, minimum height=1.1cm,
    inner sep=5pt, align=center, font=\small\bfseries
  },
]

%% ===== 节点行（y=0），水平间距 2.8cm =====
\node[datanode]    (x)  at ( 0.0, 0) {$\mathbf{x}$\\输入};
\node[opnode=blue] (W1) at ( 2.8, 0) {$W_1$\\线性层};
\node[datanode]    (a1) at ( 5.6, 0) {$\mathbf{a}_1$\\$=W_1\mathbf{x}$};
\node[opnode=green](sg) at ( 8.4, 0) {$\sigma(\cdot)$\\激活};
\node[datanode]    (h1) at (11.2, 0) {$\mathbf{h}_1$\\$=\sigma(\mathbf{a}_1)$};
\node[opnode=blue] (W2) at (14.0, 0) {$W_2$\\线性层};
\node[datanode]    (a2) at (16.8, 0) {$\mathbf{a}_2$\\$=W_2\mathbf{h}_1$};
\node[lossnode]    (L)  at (19.6, 0) {$\mathcal{L}$\\损失};

%% ===== 前向传播（蓝色实线，上方）=====
\foreach \from/\to in {x/W1, W1/a1, a1/sg, sg/h1, h1/W2, W2/a2, a2/L}{
  \draw[->, very thick, blue!70!black] (\from) -- (\to);
}
% 前向标签
\node[font=\scriptsize, blue!70!black] at (1.4, 0.75) {前向};

%% ===== 反向传播梯度（红色虚线，走节点下方深弧）=====
% 每段独立走弧，避免穿越节点。from=右节点.south, to=左节点.south
% 弧深度通过 out/in 角度控制

% L -> a2
\draw[->, very thick, red!70!black, dashed,
      out=-120, in=-60]
  (L.south) to node[below, font=\scriptsize, red!80!black, yshift=-2pt]
  {$\dfrac{\partial\mathcal{L}}{\partial\mathbf{a}_2}$}
  (a2.south);

% a2 -> h1 (跨 W2)
\draw[->, very thick, red!70!black, dashed,
      out=-130, in=-50]
  (a2.south) to node[below, font=\scriptsize, red!80!black, yshift=-2pt]
  {$\dfrac{\partial\mathcal{L}}{\partial\mathbf{h}_1}$}
  (h1.south);

% h1 -> a1 (跨 sigma)
\draw[->, very thick, red!70!black, dashed,
      out=-130, in=-50]
  (h1.south) to node[below, font=\scriptsize, red!80!black, yshift=-2pt]
  {$\dfrac{\partial\mathcal{L}}{\partial\mathbf{a}_1}$}
  (a1.south);

% a1 -> x (跨 W1)
\draw[->, very thick, red!70!black, dashed,
      out=-130, in=-50]
  (a1.south) to node[below, font=\scriptsize, red!80!black, yshift=-2pt]
  {$\dfrac{\partial\mathcal{L}}{\partial\mathbf{x}}$}
  (x.south);

% 反向标签
\node[font=\scriptsize, red!80!black] at (18.4, -2.1) {反向};

%% ===== 参数梯度（向上箭头，离开节点顶部）=====
% ∂L/∂W1
\draw[->, thick, red!60!black, dashed]
  (W1.north) -- ++(0, 1.4)
  node[above, font=\scriptsize, red!80!black, align=center]
  {$\dfrac{\partial\mathcal{L}}{\partial W_1}$};

% ∂L/∂W2
\draw[->, thick, red!60!black, dashed]
  (W2.north) -- ++(0, 1.4)
  node[above, font=\scriptsize, red!80!black, align=center]
  {$\dfrac{\partial\mathcal{L}}{\partial W_2}$};

%% ===== 图例框（右上角）=====
\node[font=\small, fill=white, draw=gray!50, thin, rounded corners=4pt,
      inner sep=7pt, align=left]
  at (15.0, 2.8) {
    \textcolor{blue!70!black}{$\longrightarrow$}\ 前向传播（Forward Pass）\\[3pt]
    \textcolor{red!70!black}{$\dashrightarrow$}\ 反向梯度（Backward Pass）
  };

%% ===== 底部公式框 =====
\node[font=\normalsize, fill=white, draw=gray!50, thin, rounded corners=5pt,
      inner sep=10pt, align=center]
  at (9.8, -5.0) {
    \textbf{矩阵微积分链式法则（反向传播核心）：}\\[6pt]
    $\displaystyle\frac{\partial\mathcal{L}}{\partial W_1}
     = \frac{\partial\mathcal{L}}{\partial\mathbf{h}_1}
       \cdot\frac{\partial\mathbf{h}_1}{\partial\mathbf{a}_1}
       \cdot\frac{\partial\mathbf{a}_1}{\partial W_1}
     = \boldsymbol{\delta}_1 \cdot \sigma'(\mathbf{a}_1) \cdot \mathbf{x}^T$
  };

%% ===== 整体标题 =====
\node[font=\large\bfseries] at (9.8, 2.8)
  {矩阵微积分与反向传播 DAG};

\end{tikzpicture}
\end{document}
